% !TEX root = ../main.tex
\section{Results}
\label{sec:results}
All values are on \modelSmall{}; the larger reward models are in progress. Numbers pulled from the experiment
JSONs are typeset from single-source macros and never hand-edited. Hiring values were produced on the earlier
\emph{synthetic} CV substrate, which the domain has since replaced with real biographies
(Section~\ref{sec:pipeline}); those arms have not yet been re-run. \outdated{Credit values (marked \dag) were produced before the German
Credit code-table correction and the factorial redesign (Section~\ref{sec:pipeline}) and must be re-run;
ASAP values (marked \ddag) predate a minor cleaning fix affecting 22 of 12,976 essays; values marked \S{}
predate the switch to the validated name lists of \citet{haim2024whatsinaname} or the domain-specific
decision verdicts.}
\outdated{Methodology decision 2026-09-24: the direct-scoring numbers below measure that the reward is
sensitive to protected attributes under controlled substitution (the mechanism layer), not that the model
assesses applicants in a biased way. The harm evidence is the cross-marker decision design
(Section~\ref{sec:method}), which has not been run at scale yet; this section is rewritten once it has.}

\subsection{Direct scoring (auto-influence)}
On the hiring (CV) domain, explicit-sex \autoinf{} is \autoInflSexCVbase{}, collapsing to
\autoInflSexCVnull{} when the \diffmean{} direction is nulled; the sex$\times$age$\times$family intersection
behaves likewise (\autoInflIntersectionCVbase{} $\to$ \autoInflIntersectionCVnull{}) and is only
\emph{partially additive} (cosine with the summed marginals \additivityCosCV{}), i.e.\ a genuine interaction.
Credit replicates the pattern (sex \autoInflSexCreditbase{} $\to$ \autoInflSexCreditnull{}).
\outdated{Credit sex pairs are now cut from the factorial, so age and marital status are stated on both
sides; the credit intersection is now sex $\times$ age $\times$ \emph{marital status}, not family status.} A robustness
battery reveals an encoding twist on CVs: the RM reacts more to realistic \emph{proxy} cues than to bald
statements for age and family status (age \batteryAgeExplicitCV{} explicit vs \batteryAgeProxyCV{} proxy;
family \batteryFamilyExplicitCV{} vs \batteryFamilyProxyCV{}). All effects are behaviourally low-complexity:
the $\alpha$-sweep drives them to near zero.

\subsection{Cross-influence (reliability harm)}
\outdated{Retired in its direct form (2026-09-24). It compared a strong and a weak record recited as the
assistant turn, a premise that did not hold: the model was grading an off-task recitation, and on unmatched
pairs a scorer preferring the longer text reached 0.97 on PERSUADE and 0.74 on ASAP. The earlier readings
(acc\_baseline 0.51 credit, 0.59 CV, 0.72 ASAP, 0.88 PERSUADE, on the MLX prototype) are not comparable and are
not reported. Decision-format \crossinf{} (Section~\ref{sec:method}) awaits the cross-marker runs.}

\subsection{Which biases, across domains}
The education domain replicates the direct-scoring picture on real essays: stating sex is a large bias
(auto-influence \eduSexExplicitAsap{} on ASAP, \eduSexExplicitPersuade{} on PERSUADE), and an education-stage
proxy is the largest single effect (\eduGradeProxyAsap{} / \eduGradeProxyPersuade{} --- the RM rewards the
identical essay far more when it is attributed to a university student rather than a school pupil), while
name-based ethnicity is weak and template-unstable (explicit \eduEthnicityExplicitAsap{} on ASAP vs
\eduEthnicityExplicitPersuade{} on PERSUADE).

\subsection{Reasoning vs.\ fairness}
Modelling the RM's scored \emph{response} rather than the applicant, the RM rewards factually correct causal
reasoning (\reasonCorrectnessPL{} on a parental-leave premise) essentially equally on a non-demographic
commute control (\reasonCorrectnessCommute{}): a domain-general reasoning-quality effect, not demographic
fairness. Consistently, the decision-response discriminatory-win-rate is near zero for sex
(\decisionDiscWinSex{}) and age (\decisionDiscWinAge{}) but rises on the legally salient family and
intersectional axes (\decisionDiscWinFamily{}, \decisionDiscWinIntersection{}).
\outdated{The decision-response arm now covers all three domains, each with its own binary decision (hiring:
advance to interview; credit: approve the loan; education: pass or fail), draws only strong applicants, and
no longer reads the synthetic ``years of experience'' field; these hiring values predate that. Since
2026-09-24 this arm is the floor --- does the model at least penalise openly stated discrimination? --- not the
harm evidence: its verdict states the reason outright and shows only the protected pole, so a low win-rate
says nothing about silent disparate treatment, which the cross-marker design measures.}

\subsection{Standpoint-credibility (positioned arguments)}
When the same argument is grounded in a stated first-person identity, the RM rewards it more from a
marginalized standpoint: on PERSUADE the intersectional identity gap (reward of the marked vs.\ reference
pole, relative to a no-standpoint baseline) is \standpointGapIntersectionPersuade{}, with sex
\standpointGapSexPersuade{}, race \standpointGapRacePersuade{}, class \standpointGapClassPersuade{}, and
origin \standpointGapOriginPersuade{}, while genuinely neutral non-demographic controls sit at approximately
zero (hobby \standpointGapHobbyPersuade{}, pet \standpointGapPetPersuade{}). ASAP replicates the ordering at
larger magnitude (intersection \standpointGapIntersectionAsap{}). The effect is robust across injection
position, corpus, three paraphrasings, and a neutral (non-endorsing) stance, and sits on top of a general
reward for any lived-experience appeal (main effect \standpointMainEffectPersuade{}).

\subsection{Complexity: \leace{} + non-linear probe}
The decisive test: after \leace{}, the linear probe for the reasoning-correctness concept drops to
\erasureLinearLeaceCorrPL{} (chance, by construction) while an MLP still recovers it at
\erasureMlpLeaceCorrPL{}. The concept is linearly erasable but non-linearly recoverable, the \taco{}
signature of entanglement. The reward-scalar low-complexity reading is therefore incomplete: null-space
projection collapses the bias's contribution to the reward, but the concept remains present in the
representation.
\emph{Scope.} This test has so far been run only on the reasoning concepts of the model-response arm, in the
hiring domain. The corresponding test on the \emph{demographic} directions --- which would establish the same
gap for protected attributes directly --- has not yet been run, and is our immediate next experiment.
